\section{Metodologie e Dati}\label{sec:methods}

Diversi metodi sono stati impiegati, per dimostrare risultati diversi e per
averne maggiore riscontro.\\
Per la causalità si è usata la regressione; Granger non non poteva essere
utilizzato, in mancanza di stazionarietà delle serie di GDP e medaglie (misurata
con ADF test), pure applicando modifiche come logaritmo, differenza di
logaritmi, percentuale (medaglie).\\
I periodi analizzati terminano nel 2026 (anno di Giochi Invernali) e partono in
uno tra 1996, 1964, 1932, 1900: la scelta di ognuno minimizza il rumore dovuto
agli anni appena precedenti (per esempio, guardare il 1948 significherebbe usare
PIL della seconda guerra mondiale). In particolare, il 1932 è distanziato dalla
prima guerra mondiale, il 1964 dalla seconda, il 1996 dalla guerra fredda; il
1900 rappresenta la seconda edizione (la prima va scartata per fini statistici).
I quattro anni appartengono anche a periodi diversi per numero di partecipanti,
rispettivamente 25, 41, 85, 164.\\

\subsection{Dati}\label{sec:methods_data}



\section{Regressione}
\label{sec:methods}

Il nostro lavoro non solo espande quello di quello di Csurilla\cite{csurillaLessObviousEffect2023} analizzando altri fattori, 
ma impiega anche un modello regressione diverso.\\
Non avendo effettuato un'analisi a livello di sport, infatti, cercare di riprodurre fedelmente i suoi risultati porta ad 
erorri di esecuzione della regressione, con molta probabilità dovuti a casi di collinearità (perfetta o quasi) tra variabili, 
a cui ZINB è molto sensibile. Sempre per evitare collinearità perfetta, quando si usa \textit{YEAR}, bisogna escludere un anno:
il suo effetto è già assorbito da \textit{CONST}.\\
\todo{espandere, anche con mie note}